\documentclass[runningheads]{ceurart}
\usepackage{graphicx}
\usepackage{booktabs}
\usepackage{amsmath}
\usepackage{microtype}
\usepackage{caption}
\captionsetup{font=small,skip=4pt}
\usepackage{hyperref}
\usepackage{multirow}
\usepackage{float}

\usepackage[T1]{fontenc}
\usepackage[utf8]{inputenc}
\usepackage{placeins}
\usepackage{enumitem}
\setlist{nosep}
% Reduce space around floats and lists for compactness
\setlength{\textfloatsep}{8pt}
\setlength{\floatsep}{8pt}
\setlength{\intextsep}{8pt}
\setlength{\parskip}{0pt}
\setlength{\itemsep}{0pt}

\begin{document}

\copyrightyear{2025}
\copyrightclause{Copyright for this paper by its authors. Use permitted under Creative Commons License Attribution 4.0 International (CC BY 4.0).}
\conference{Forum for Information Retrieval Evaluation, December 2025, India}

\title{Enhanced Code Comment Relevance Classification Using Domain-Adaptive CodeBERT with Masked Language Modeling}

\author[1]{L. Simiyon Vinscent Samuel}
\author[1]{Sundareswaran Rajaguru}
\author[1]{Vijayaraaghavan K. S.}
\author[1]{A. Vijayalakshmi}

\address[1]{Department of Computer Science and Engineering,
Sri Sivasubramaniya Nadar College of Engineering,
Chennai,India}



\begin{abstract}
We evaluate domain-adaptive pre-training of CodeBERT using masked language modeling (MLM) for binary code-comment relevance classification in the IRSE shared task. We compare CodeBERT+MLM against TF--IDF baselines and study the effect of synthetic and silver-standard augmentation. CodeBERT+MLM attains F$_1$=0.9193 on the original test split versus 0.7067 for the best TF--IDF baseline (logistic regression). Augmentation (combined data) produces a marginal change (F$_1$=0.9166), highlighting that augmentation quality — not quantity — drives gains. We provide implementation details and practical recommendations for applying domain-adaptive pre-training in software engineering tasks.
\end{abstract}

\begin{keywords}
code comment classification \sep CodeBERT \sep domain-adaptive pre-training \sep masked language modeling \sep data augmentation
\end{keywords}

\maketitle

\section{Introduction}
Code comments are essential documentation artifacts that improve comprehension and maintenance~\cite{rani2023decade,majumdar2022automated}. The IRSE shared task frames comment relevance as a binary classification problem: whether a comment meaningfully describes a code segment. While classical TF--IDF models often perform strongly on technical text~\cite{basu2023identification}, pre-trained transformers tailored to code (e.g., CodeBERT~\cite{feng2020codebert}) can capture richer cross-modal patterns. Domain-adaptive pre-training (DAPT) via continued MLM on domain text has been shown to close distributional gaps between generic pretraining and task data~\cite{gururangan2020dapt}. 

This paper investigates CodeBERT with intermediate MLM on C code-comment corpora and compares it to TF--IDF baselines, assessing the impact of rule-based synthetic and silver-standard augmentations. Our main contributions: (i) an empirical evaluation of CodeBERT+MLM on original and augmented data; (ii) comparisons with classical baselines; and (iii) analysis of augmentation effects and cross-dataset generalization.

\section{Related work (condensed)}
Transformers adapted for code, such as CodeBERT, enable joint code–text representations useful for search and summarization~\cite{feng2020codebert,liu2022exploiting}. DAPT/TAPT improves downstream task fit by continued MLM on domain/task data~\cite{gururangan2020dapt,kim2025subset}. Prior IRSE work shows TF--IDF models can be competitive for comment relevance, motivating careful empirical comparison~\cite{basu2023identification}. Data augmentation (rule-based or model-based) helps low-resource settings but can introduce label noise in domain-specific tasks lacking automatic verification~\cite{wang2023synthetic,feng2021augmentation}.

\section{Datasets and experimental setup}
\subsection{Corpora and silver-standard creation}
We use the IRSE/FIRE seed dataset (11,452 code–comment pairs extracted from C projects). After cleaning malformed or unlabeled entries we obtain 8,326 training samples (``Original''). To expand data we used two augmentation sources:

\begin{enumerate}
  \item \textbf{Rule-based synthetic:} 500 examples generated by templated paraphrasing (synonym substitution, sentence restructuring) derived from existing comments to increase stylistic diversity while preserving labels.
  \item \textbf{HumanEval-derived silver standard:} 1,640 examples extracted from \texttt{human-eval.jsonl}. We parsed each example to extract function docstrings and surrounding comments (mostly Python), then used an OpenRouter ``fast-grok'' free model as a preprocessing / normalization step to:
    \begin{itemize}
      \item split function-level docstrings into candidate comment sentences,
      \item canonicalize parameter and return descriptions,
      \item filter out non-descriptive or placeholder docstrings (e.g., single-word comments).
    \end{itemize}
  The normalized candidates were programmatically labeled using heuristics (presence of function-level description, param/return mentions, and overlap with code identifiers). These heuristics produce silver-standard labels (higher noise than human annotation but useful signal).
\end{enumerate}

We then combined the C-based Original training set with Python-derived silver-standard and rule-based synthetic examples to form the \textbf{Combined} training set (8,826 samples). Table~\ref{tab:data-stats} summarizes composition.

\begin{table}[h]
    \centering
    \caption{Dataset composition after cleaning.}
    \label{tab:data-stats}
    \small
    \begin{tabular}{lrrr}
        \toprule
        Dataset & Samples & Useful & Not Useful \\
        \midrule
        Original seed & 11,452 & 7,063 (61.7\%) & 4,389 (38.3\%) \\
        Rule-based synthetic & 500 & 235 (47.0\%) & 265 (53.0\%) \\
        HumanEval silver & 1,640 & 970 (59.1\%) & 670 (40.9\%) \\
        \midrule
        Training (Original) & 8,326 & -- & -- \\
        Training (Combined) & 8,826 & -- & -- \\
        \bottomrule
    \end{tabular}
\end{table}

\subsection{Preprocessing and labeling notes}
Code and comment texts were concatenated with a separator token. Tokenization uses CodeBERT's fast BPE tokenizer with a maximum sequence length of 512. For silver-standard examples we applied two lightweight filters: (1) minimum token length (≥3 words) and (2) presence of at least one verb or function-related identifier, to reduce trivial or metadata-only docstrings. The silver-standard labeling heuristics prioritize precision (favoring conservative "useful" labels) but inevitably inject label noise; we therefore treat these examples as complementary rather than authoritative.

\subsection{Models and training}
We continued MLM from \texttt{microsoft/codebert-base-mlm} on the Original (C) training corpus prior to supervised fine-tuning (15\% mask, batch size 16, lr 5e-5, max 7 epochs, AdamW). After DAPT we replace the MLM head with a classification layer (dropout 0.1) and fine-tune on the binary task (batch size 32, lr 2e-5, up to 10 epochs, linear warmup 10\% steps). Baselines use TF--IDF features with: logistic regression (L2), linear SVM, multinomial NB, and random forest (100 trees). Baselines use 10-fold CV for hyperparameter tuning.

\subsection{Evaluation}
Metrics: accuracy, precision, recall, weighted F$_1$. Experimental configurations: (A) train Original → test Original; (B) train Combined → test Combined; (C) Original → Combined; (D) Combined → Original. Results are reported on held-out test splits not used during validation.

\subsection{Why MLM trained on C still helps after mixing Python silver data}
Although silver examples derive from Python docstrings, the MLM DAPT step was performed primarily on C code and comments (Original corpus). Two reasons explain why DAPT remained effective after adding Python silver examples for classification:
\begin{enumerate}
  \item \textbf{Cross-language lexical overlap:} many identifier names, API-style phrases, and comment conventions (e.g., ``returns'', ``params'', ``compute'') are shared across languages; the C-trained MLM learns these patterns and applies them when encoding Python-derived comments.
  \item \textbf{Representation robustness:} MLM adapts CodeBERT's token and contextual embeddings to code/comment distributions (punctuation, identifier tokens, inline code snippets). These representational improvements transfer to mixed-language fine-tuning because the downstream classifier primarily relies on semantic patterns rather than language-specific syntax.
\end{enumerate}
Empirically, mixing Python silver data produced only a slight drop in peak F$_1$ (0.9193 $\rightarrow$ 0.9166), indicating that distributional mismatch is limited and quality/noise in silver labels is the larger contributor to performance change.

\section{Results}
\subsection{CodeBERT+MLM performance}
Table~\ref{tab:transformer} reports CodeBERT+MLM across configurations. The model achieves peak in-distribution F$_1$=0.9193 (Part A). Augmentation produces a small change (Part B F$_1$=0.9166).

\begin{table}[h]
    \centering
    \caption{CodeBERT+MLM performance (held-out tests).}
    \label{tab:transformer}
    \small
    \begin{tabular}{lccc}
        \toprule
        Configuration & Accuracy & F$_1$ & Precision \\
        \midrule
        Original $\rightarrow$ Original & 0.9196 & 0.9193 & 0.9193 \\
        Combined $\rightarrow$ Combined & 0.9179 & 0.9166 & 0.9193 \\
        Original $\rightarrow$ Combined & 0.9145 & 0.9141 & 0.9141 \\
        Combined $\rightarrow$ Original & 0.9190 & 0.9176 & 0.9203 \\
        \bottomrule
    \end{tabular}
\end{table}

\begin{figure}[h]
    \centering
    \includegraphics[width=0.85\textwidth]{codebert_ml_comparison.png}
    \caption{Performance comparison overview: CodeBERT+MLM vs TF--IDF baselines.}
    \label{fig:codebert_comparison}
\end{figure}


\subsection{Baselines and dynamics}
Table~\ref{tab:baselines} shows TF--IDF baseline results. Logistic regression is the strongest baseline (F$_1$=0.7067 on Combined), well below CodeBERT+MLM. Figure~\ref{fig:training_progress} shows training dynamics for MLM and fine-tuning.

\begin{table}[h]
    \centering
    \caption{TF--IDF baselines (accuracy / F$_1$).}
    \label{tab:baselines}
    \small
    \begin{tabular}{llcc}
        \toprule
        Model & Dataset & Accuracy & F$_1$ \\
        \midrule
        Logistic Regression & Original & 0.7203 & 0.6897 \\
         & Combined & 0.7288 & 0.7067 \\
        SVM & Original & 0.7041 & 0.6796 \\
         & Combined & 0.7163 & 0.6985 \\
        Naive Bayes & Original & 0.7023 & 0.6676 \\
         & Combined & 0.7123 & 0.6925 \\
        Random Forest & Original & 0.6771 & 0.6644 \\
         & Combined & 0.7050 & 0.6999 \\
        \bottomrule
    \end{tabular}
\end{table}

\begin{figure}[!htbp]
    \centering
    \includegraphics[width=0.85\textwidth]{task1_training_progress.png}
    \caption{Training/validation trends (MLM + fine-tune): rapid MLM perplexity reduction and stable fine-tuning convergences.}
    \label{fig:training_progress}
\end{figure}

MLM reduced perplexity notably over 7 epochs (from 15.3 to 8.7). Fine-tuning converged in 6–8 epochs with early stopping. Augmentation (+6\% data) produced a 0.27pp F$_1$ decrease; manual inspection indicates silver-standard noise and stylistic mismatch as main causes.

\subsection{Comparison with previous IRSE / FIRE submissions}
We compare our best performance against published IRSE-relevant works (table below uses provided summary metadata). This contextualizes gains from DAPT+MLM.

\begin{table}[h]
    \centering
    \caption{Selected previous IRSE/FIRE submissions (summary).}
    \label{tab:prior_irse}
    \small
    \begin{tabular}{lllc}
        \toprule
        Paper ID & Title (short) & Primary model(s) & Best reported F$_1$ \\
        \midrule
        T7-1 & Generative AI for Software Metadata (Majumdar et al.) & NN, SVM, RF & 0.885 \\
        T7-2 & ML-LLM pairing (Abi Akl) & NN, RF, Voting & 0.884 \\
        T7-3 & Software Metadata + GenAI & SVM, ANN & 0.85 (approx) \\
        T7-6 & Enhancing Binary Comment Quality & BERT, DT, ANN & 0.885 \\
        T7-9 & Leveraging LMs for Code Comment & GPT-3 / Codex & 0.96 (LLM results) \\
        \midrule
        This work & CodeBERT + DAPT (MLM) & CodeBERT+MLM & \textbf{0.9193} \\
        \bottomrule
    \end{tabular}
\end{table}

Our CodeBERT+MLM result (F$_1$=0.9193) exceeds many classical/ML+LLM-augmented submissions while remaining below reported very large LLM one-off results (T7-9/T7-10) that use generative LMs as primary models. The comparison suggests DAPT is a practical and effective middle ground: much better than TF--IDF baselines and less resource-intensive than training or prompting very large LLMs end-to-end.

\begin{figure}[!htbp]
    \centering
    \includegraphics[width=0.85\textwidth]{performance_comparison.png}
    \caption{Final metric comparison (CodeBERT+MLM vs baselines).}
    \label{fig:perf_comparison}
\end{figure}
\FloatBarrier  
\section{Discussion}
Domain-adaptive MLM significantly improves CodeBERT representations for comment relevance. Cross-evaluation shows training on Original and testing on Combined drops F$_1$ by only 0.52pp, indicating robust transfer. The small degradation with augmentation confirms that label quality and distribution alignment matter: silver-standard Python-derived examples add vocabulary and diversity, but their programmatic labels and stylistic differences introduce noise that slightly offsets gains.

Limitations: single-language (C) focus in DAPT, binary relevance label space, and the computational cost for continued pretraining. Future work: multi-task DAPT, curriculum strategies that gradually introduce silver examples, active human verification of synthetic examples, and interpretability analyses to surface tokens driving decisions.

\section{Conclusion}
We show CodeBERT with domain-adaptive MLM attains strong results on comment relevance classification, outperforming TF--IDF baselines by a wide margin. Silver-standard augmentation from HumanEval (processed via OpenRouter fast-grok + heuristics) provides additional data but must be quality controlled to benefit high-capacity transformers. DAPT is a practical, resource-efficient strategy to improve specialized SE NLP tasks.

\begin{thebibliography}{99}

\bibitem{rani2023decade}
Rani, P., Panichella, S., Leuenberger, M., Ghafari, M., Nierstrasz, O.: 
A decade of code comment quality assessment: A systematic literature review. 
J. Syst. Softw. 196, 111515 (2023). 
\url{https://doi.org/10.1016/j.jss.2022.111515}

\bibitem{majumdar2022automated}
Majumdar, S., Liang, Y., Ripley, M., Mody, N., Dalal, K., Chen, L.: 
Automated evaluation of comments to aid software maintenance. 
Empir. Softw. Eng. 27(5) (2022). 
\url{https://doi.org/10.1007/s10664-022-10171-5}

\bibitem{fire2024}
Forum for Information Retrieval Evaluation (FIRE). 
\url{https://fire.irsi.org.in} (2024)

\bibitem{irse2023overview}
Majumdar, S., Paul, S., Paul, D., Bandyopadhyay, A., Chattopadhyay, S., Das, P.P., Clough, P.D., Majumder, P.: 
Generative AI for software metadata: Overview of the Information Retrieval in Software Engineering Track at FIRE 2023. 
arXiv:2311.03374 (2023). 
\url{https://arxiv.org/abs/2311.03374}

\bibitem{feng2020codebert}
Feng, Z., Guo, D., Tang, D., Duan, N., Feng, X., Gong, M., Shou, L., Qin, B., Liu, T., Jiang, D., Zhou, M.: 
CodeBERT: A pre-trained model for programming and natural languages. 
In: Findings of EMNLP, pp. 1536--1547 (2020). 
\url{https://doi.org/10.18653/v1/2020.findings-emnlp.139}

\bibitem{basu2023identification}
Basu, T., Sruthi, S.: 
Identification of the relevance of comments in codes using bag of words and transformer based models. 
arXiv:2308.06144 (2023). 
\url{https://arxiv.org/abs/2308.06144}

\bibitem{gururangan2020dapt}
Gururangan, S., Marasovi\'{c}, A., Swayamdipta, S., Lo, K., Beltagy, I., Downey, D., Smith, N.A.: 
Don't stop pretraining: Adapt language models to domains and tasks. 
In: Proceedings of ACL, pp. 8342--8360 (2020). 
\url{https://doi.org/10.18653/v1/2020.acl-main.740}

\bibitem{liu2022exploiting}
Liu, K., Kim, D., Bissyand\'{e}, T.F., Yoo, S., Le Traon, Y.: 
Better exploiting CodeBERT to support source code-related tasks. 
In: ACM International Conference on Mining Software Repositories (MSR), pp. 545--549 (2022). 
\url{https://doi.org/10.1145/3545258.3545260}

\bibitem{mohammadkhani2022explaining}
Mohammadkhani, A.H., Panichella, S., Parsa, S.: 
Explaining transformer-based code models: How do they learn code syntax and semantics? 
In: IEEE International Conference on Software Maintenance and Evolution (ICSME), pp. 519--523 (2022). 
\url{https://doi.org/10.1109/ICSME58945.2022.00058}

\bibitem{chen2024prompt}
Chen, X., et al.: 
Enhancing source code classification effectiveness via prompt engineering. 
Nat. Sci. Rep. 14, 69402 (2024). 
\url{https://doi.org/10.1038/s41598-024-69402-7}

\bibitem{kim2025subset}
Kim, S., Lee, J., Park, M.: 
Subset selection for domain adaptive pre-training of language model. 
Nat. Sci. Rep. 15, 94085 (2025). 
\url{https://doi.org/10.1038/s41598-025-94085-z}

\bibitem{wang2023synthetic}
Wang, Y., Zheng, Q., Chen, H., Yang, B., Xu, J., Gao, M.: 
Synthetic data generation with large language models for text classification. 
arXiv:2310.11467 (2023). 
\url{https://arxiv.org/abs/2310.11467}

\bibitem{feng2021augmentation}
Feng, S.Y., et al.: 
A survey of data augmentation approaches for NLP. 
In: ACL Findings, pp. 968--988 (2021). 
\url{https://doi.org/10.18653/v1/2021.findings-acl.84}

% IRSE 2023 Track Papers from CEUR-WS Vol-3681
\bibitem{T7-1}
Majumdar, S., et~al.: 
Overview of the Information Retrieval in Software Engineering Track at FIRE 2023. 
In: CEUR Workshop Proceedings, vol. 3681 (2023). 
\url{https://ceur-ws.org/Vol-3681/T7-1.pdf}

\bibitem{T7-6}
Arumugam, R., Deborah, A.: 
Enhancing binary code comment quality classification: Integrating generative AI for improved accuracy. 
In: CEUR Workshop Proceedings, vol. 3681 (2023). 
\url{https://ceur-ws.org/Vol-3681/T7-6.pdf}

\bibitem{T7-7}
Binary classification of source code comments using machine learning. 
In: CEUR Workshop Proceedings, vol. 3681 (2023). 
\url{https://ceur-ws.org/Vol-3681/T7-7.pdf}

\end{thebibliography}

\end{document}
